% New chapter for Part I of BODYandSOUL: MATHEMATICAL SIMULATION TECHNOLOGY.
% Intended to follow Chapter 2, "From Pythagoras to Google".  Draft 2026.
% NOTE: cross-references are written as plain chapter numbers from the 2010 draft;
% they should become \chref{} / \sref{} once the Part I source is available.

\chapter{From Google to GPT}

\begin{quote}
The opposite of a truth of reasoning is impossible; the opposite of a truth of fact is possible.
\hfill (Leibniz)
\end{quote}

\section{What Has Happened}

Chapter 2 called the previous chapter \emph{From Pythagoras to Google} and section 2.1 called what
was coming \emph{The Singularity}. It arrived, and it is worth being exact about what it did.

Machines now carry out symbolic derivation. They write computer programs. They produce proofs that
machine proof-checkers accept. Anything in mathematics that consists of following a rule correctly
is now done faster and more cheaply by a machine than by a person, and this includes the activity
this book asks the student to take up: writing the code and letting the machine run it.

So the remedy proposed in the Preface has itself been automated. This chapter is about what that
leaves.

\section{Formulate and Solve}

Mathematics is two activities:

\begin{itemize}
\item[(i)] \emph{formulate} the equation;
\item[(ii)] \emph{solve} the equation.
\end{itemize}

That division is the organising structure of this book: Parts II and VIII set up equations of
motion, Parts IV, V and IX solve them, and Parts III, VII and X let the student watch the solving
happen.

Everything that has been automated belongs to (ii). Nothing has been automated in (i), and nothing
about the automation suggests it will be. There is no rule for deciding which equation describes a
situation, because that decision answers to the world and not to a rule. A curriculum organised
around the automated half is in difficulty whatever anybody decides about it; a curriculum organised
around the other half has just become the only one worth running.

\section{Is Solving Then Empty?}

There is an objection which, if it held, would make all this of no consequence.

Formal mathematics draws logical consequences from axioms. A consequence is already contained in
the axioms --- that is what makes it necessary, and it is Leibniz's own point that the opposite of a
truth of reasoning is impossible. So nothing is ever added: the axioms are posited, the theorems
unpacked, and the unpacking introduces no content the positing did not carry. On this reading
solving is a tautology of record, and automating it automates nothing.

The objection confuses what is \emph{implied} with what is \emph{available}. A chain of steps brings
out what is in no useful sense present at the start --- not an addition to the axioms, but a
\emph{construction} built upon them. The content appears in the chain, because the chain is where
the work is, and the length of the chain is not bounded by the length of the axioms.

\section{The Equation Is the Axiom Set}

This is clearest where this book already works. Solve $f(x)=0$ by Newton's method
(Chapter 75):
\[
x_{k+1} \;=\; x_k - \frac{f(x_k)}{f'(x_k)} .
\]
The equation plays the role of the axioms: it posits, and it determines its solution without
exhibiting it. The solution plays the role of the theorem: implicit from the outset, unavailable
until constructed. The iteration plays the role of the proof: a chain of legitimate steps which
makes available what was only implied.

Nobody is tempted to call Newton's method empty. It is the same structure, seen without
controversy, and it shows what the objection misses: the positing and the constructing are
different acts, and only the second has been handed to machines.

\section{The Two Kinds of Truth}

Leibniz divided truths of \emph{reasoning} from truths of \emph{fact}. The first follow from what
is granted, and a machine manipulating symbols correctly can produce them all. The second must be
checked against how things are, and no manipulation produces them.

What has been automated is exactly the first and none of the second. That is why the question of
what to teach now has an answer rather than merely an anxiety. Section 2.3 of this book is called
\emph{From Formal to Real Knowledge}; the machines have made that passage compulsory.

It also settles, in an unexpected way, the quarrel the Preface describes between the
formalist/logicist school and the constructive one. Formal derivation from axioms is precisely what
machines do well. It is not wrong. It has become \emph{cheap} --- and cheapness decides curricula
where argument never did.

\section{What Is Left for the Student}

Three things, and they are not small.

\paragraph{Formulating.} Deciding what equation describes the situation. This is the whole of (i),
it is answerable to the world, and it is what the simulators of Parts III, VII and X exist to teach.

\paragraph{Judging.} Deciding whether what the machine returns is correct. A machine that writes
code writes plausible code, and plausible is not correct. Judging it requires knowing what the
method does and what it costs --- which is the Soul of this book's title and was always its reason
for existing.

\paragraph{Understanding why it works.} Not in order to do it by hand, which nobody will do again,
but in order to do the judging. A person who has only watched cannot tell a converged computation
from one that has merely stopped moving.

\section{The New Exercise}

The practical change is a reversal, and this book is unusually ready for it.

In 2010 the student wrote the simulator and the machine ran it. Now the student \emph{specifies} the
simulator, the machine writes it, and the student must decide whether the result is right. That is
the harder task, not the easier one.

The simulators in Parts III, VII and X are well suited to it precisely because their answers are
known here. Specify the pendulum, the elastic string, the planetary system to a machine; let it
produce the code; then check what comes back against the analysis in this book. \emph{Catching the
machine where it is wrong is the exercise}, and it cannot be done without the mathematics.

\section{Objections}

\paragraph{``The machines will learn to formulate too.''} Perhaps. Nothing in what they do now
suggests it, because formulating is answerable to the world and their training is answerable to
text. But the claim here does not require the limit to be permanent: it requires only that the two
halves be automated at different times, which has already happened.

\paragraph{``Then nobody need learn to compute.''} The opposite. Somebody must be able to tell a
correct computation from a confident one, and that ability has no shortcut.

\paragraph{``This is just the calculator argument again.''} The calculator automated arithmetic and
the response was to teach algebra. The machine automates algebra, and the response is to teach
modelling. The pattern is the same and so is the direction: up one level, toward the part that
answers to the world.
